% Cote 119: offprints, reports, courses, CV and letters (typed and printed part): archive cote 119, site page 65 (PDF page 66).
% offprint. Draft compared: none.
% AI-transcribed by Claude Opus 5.5 (claude-opus-5-5, adaptive thinking, effort high), Message Batches API, batch r119_full (msgbatch_01DwZNm5Y7YsKBWyAN4sFs8w), 2026-11-13; not yet checked by eye.
\scanpage{119/65}
\mspage{C. R. Acad. Sc. Paris, t. 273, p. 1246-1248 (20 décembre 1971) \quad Série A}

GÉOMÉTRIE ALGÉBRIQUE. --- Polarisations des catégories tannakiennes :
cas gradué. Note (*) de M. Neantro Saavedra Rivano, transmise par
M. Henri Cartan.

Introduction. --- On développe ici la variante graduée de la notion
de polarisation introduite dans une Note précédente (${}^{1}$). C'est sous cette
forme que cette structure se présente dans les exemples qui ont motivé
la définition et l'étude des polarisations [voir (${}^{2}$), 1.5, 5.11 et le chapitre
des exemples dans (${}^{3}$)].
On se donne une catégorie tannakienne algébrique $C$ sur $\mathbf{R}$; le lien
de $C$ sera noté L, la L-gerbe correspondant à $C$ sera notée $\mathcal{G}$ [voir (${}^{4}$)
pour les généralités concernant les catégories tannakiennes], et le centre
de L sera noté Z
\[ \mathrm{Z} = \mathit{Aut}^{\otimes}(\mathrm{id}_{C}) = \mathit{Aut}(\mathrm{id}_{\mathcal{G}}). \]

On remarque que la donnée d'une $\otimes$-graduation de type $\mathbf{Z}$ de $C$ (ou plus
précisément, du $\otimes$-foncteur $\mathrm{id}_{C}$) revient à celle d'un morphisme de
$\mathbf{R}$-groupes $w : \mathbf{G}_{m} \to \mathrm{Z}$. De façon analogue, la donnée d'un objet inver-%
sible T de $C$ revient à celle d'un morphisme de gerbes $t : \mathcal{G} \to \mathrm{TORS}(\mathrm{G}_{m})$;
elle détermine un morphisme, noté encore $t$, $t : \mathrm{Z} \to \mathbf{G}_{m}$. Si on dispose
de deux telles données $w$, T, le poids de T pour la graduation $w$ est
l'entier $n$ tel que $t \circ w(\lambda) = \lambda^{n}$ pour $\lambda \in \mathbf{G}_{m}$.

1. Triples de Tate.

Définition 1.1. --- Un triple de Tate $T$ est un triple $T = (C, w, \mathrm{T})$
constitué d'une catégorie tannakienne $C$ sur $\mathbf{R}$, une $\otimes$-graduation
$w : \mathbf{G}_{m} \to \mathrm{Z}$ de type $\mathbf{Z}$ de $\mathrm{id}_{C}$ et d'un objet inversible T de $C$ qui soit
de poids $-2$.

1.2. Si $T$ est un triple de Tate, notons $\mathcal{G}_{0}$ la gerbe noyau du morphisme
$t : \mathcal{G} \to \mathrm{TORS}(\mathrm{G}_{m})$; ses objets sont les couples $(\mathrm{P}, \xi)$ d'un objet P
de $\mathcal{G}$ et d'une trivialisation $\xi$ du torseur $t(\mathrm{P})$ sous $\mathrm{G}_{m}$
\[ \xi : \ \mathrm{P} \xrightarrow{\sim} \mathbf{G}_{m}. \]

La gerbe $\mathcal{G}_{0}$ munie de $\mathcal{G}_{0} \to \mathcal{G}$, correspond à une catégorie tannakienne $C_{0}$
munie d'un morphisme de catégories tannakiennes
\[ \mathrm{Q} : \ C \to C_{0}; \]

l'image de T par ce morphisme est un objet unité de $C_{0}$ et ce morphisme
est universel pour cette propriété. Si $\mathrm{Z}_{0}$ est le centre de $C_{0}$
\[ \mathrm{Z}_{0} = \mathit{Aut}^{\otimes}(\mathrm{id}_{C_{0}}) = \mathit{Aut}(\mathrm{id}_{\mathcal{G}_{0}}), \]

\mspage{\hand{65}}
